% Figure 1 tikzpicture — Spanish edition (same anti-collision
% geometry and sovereign nomenclature as the English version, per
% GUIA_BUENAS_PRACTICAS_LATEX_IEEE; palette seni-* from the host preamble).
\begin{tikzpicture}[
  fase/.style={rectangle, rounded corners=5pt, draw, thick, text width=11.6cm, align=center, inner sep=5pt, minimum height=9mm},
  compuerta/.style={rectangle, rounded corners=2pt, draw=seni-dark!50, thick, fill=white, text=seni-dark, font=\tiny\bfseries, minimum width=6.8mm, minimum height=4.8mm, inner sep=0pt},
  forense/.style={rectangle, rounded corners=3pt, draw=seni-red, thick, fill=seni-red!8, text=seni-dark, font=\tiny, align=center, text width=2.6cm, inner sep=2pt},
  flujo/.style={-{Stealth[length=3mm]}, very thick, seni-dark!70},
  retorno/.style={-{Stealth[length=2.5mm]}, thick, seni-red, dashed},
  nota/.style={font=\scriptsize, text=seni-dark!85, align=center, fill=white, inner sep=2pt},
]
%% --- FASE 1 ---
\node[fase, fill=seni-blue!88, text=white] (f1) at (0, 0) {
  \textbf{FASE 1 $\cdot$ ENCLAVE HERMÉTICO MACRO-SANDBOX}\\[1pt]
  {\scriptsize Espacios de nombres POSIX (clase bubblewrap) $\cdot$ FS solo-lectura $\cdot$ presupuesto ${<}\,200$ MiB}
};
\node[nota, text=seni-blue] at (0, -1.45) {\textbf{Cascada de 14 compuertas} $G_0$--$G_{13}$ $\cdot$ semántica \emph{fail-closed}};
%% --- Compuertas G0..G13 (fila única) ---
\foreach \i in {0,...,13}{
  \node[compuerta] (g\i) at (-4.55+\i*0.70, -2.15) {$G_{\i}$};
}
%% --- FASE 2 ---
\node[fase, fill=seni-green!85, text=white] (f2) at (0, -4.00) {
  \textbf{FASE 2 $\cdot$ BANCO DE INVESTIGACIÓN} (Modelo de Referencia Afinado)\\[1pt]
  {\scriptsize Verificación estática post-evento $N-1$ (catálogo de 215 contingencias) $\cdot$ telemetría del operador $\cdot$ reconciliación entre herramientas}
};
%% --- FASE 3 ---
\node[fase, fill=seni-gold!90, text=white] (f3) at (0, -6.60) {
  \textbf{FASE 3 $\cdot$ GOLDEN MASTER}\\[1pt]
  {\scriptsize Re-verificación \emph{fail-closed} de la cascada completa $\cdot$ sellado SHA-512 desacoplado por artefacto}
};
%% --- FASE 4 ---
\node[fase, fill=seni-red!88, text=white] (f4) at (0, -9.20) {
  \textbf{FASE 4 $\cdot$ PRODUCCIÓN INDUSTRIAL} (Motor de Ejecución de Red a Escala Plena)\\[1pt]
  {\scriptsize Conmutación atómica de kernel \texttt{renameat2(RENAME\_EXCHANGE)} $\cdot$ delta de regresión $= 0\%$}
};
%% --- Flujo principal (rótulos jamás cruzados: relleno blanco midway) ---
\draw[flujo] (f1.south) -- (0, -1.77);
\draw[flujo] (0, -2.53) -- node[nota, right=2mm] {14/14 PASS} (f2.north);
\draw[flujo] (f2.south) -- node[nota, right=2mm] {paridad total} (f3.north);
\draw[flujo] (f3.south) -- node[nota, right=2mm] {estado inmutable} (f4.north);
%% --- Lazo FAIL (izquierda): informe forense y retorno a la Fase 1 ---
\node[forense] (forense) at (-8.2, -2.15) {FAIL en cualquier compuerta: informe forense de causa};
\draw[retorno] (-5.14, -2.15) -- (forense.east);
\draw[retorno] (forense.north) -- (-8.2, 0) -- (f1.west);
%% --- Lazo de desviación (derecha): Fase 2 de vuelta a la Fase 1 ---
\draw[retorno] (f2.east) -- (7.2, -4.00) -- (7.2, 0) -- (f1.east);
\node[nota, text=seni-red, rotate=90, align=center] at (7.58, -2.00) {desviación no explicada $\to$ retorno a la Fase 1};
\end{tikzpicture}
